\section{Objetivos}
\label{sec:objetivos}

\subsection{Criterio de redacción}

Conforme a la indicación del profesor guía, cada objetivo explicita \textbf{qué se hará}, \textbf{cómo
se hará} y \textbf{qué resultado verificable se obtendrá}, acompañado de un indicador medible y un
plazo. La estructura aplicada es:

\begin{cajadefinicion}
\emph{[verbo en infinitivo]} + \emph{[objeto]} + \textbf{mediante} \emph{[método]} + \textbf{para
obtener} \emph{[resultado verificable]} + \textbf{en} \emph{[plazo]}.
\end{cajadefinicion}

\subsection{Objetivo general}

\begin{cajaclave}[title={Objetivo general}]
Diseñar, implementar y validar una plataforma web abierta para la gestión de datos de redes de sensores
IoT heterogéneas, que unifique la ingesta mediante un protocolo de comunicación basado en el estándar
Contact ID (SIA DC-05), aplicando la metodología PPDIOO, de modo que una organización que ya posee
sensores de distintas tecnologías pueda centralizar, normalizar y administrar sus datos desde una única
interfaz, validada sobre un caso de uso con sensores reales al término del semestre académico 2026.
\end{cajaclave}

\subsection{Objetivos específicos}

\subsubsection{OE1 --- Caracterizar el problema y establecer los requerimientos}

\begin{cajadefinicion}
\textbf{Caracterizar} la situación actual de gestión de datos de sensores en la organización
destinataria, \textbf{mediante} entrevistas semiestructuradas a las partes interesadas, revisión de
literatura técnica y análisis de la oferta regional, \textbf{para obtener} un catálogo de requerimientos
funcionales y no funcionales íntegramente trazables a las causas del problema, \textbf{al 28 de agosto
de 2026}.
\end{cajadefinicion}

\begin{longtable}{@{}L{1.5cm}L{12.2cm}@{}}
\filaencab \encab{} & \encab{Concreción} \\ \midrule \endhead
\textbf{S} & Catálogo de requerimientos con trazabilidad explícita causa $\rightarrow$ requerimiento. \\
\textbf{M} & $\geq 2$ entrevistas realizadas; $\geq 20$ requerimientos documentados; \textbf{100\%} trazados a al menos una causa (C1--C5); $\geq 2$ competidores analizados. \\
\textbf{A} & Requiere gestión de entrevistas y trabajo documental, sin dependencia de hardware. \\
\textbf{R} & Aborda C4 y sustenta los criterios 1 y 3 de la etapa, que suman el 55\% de su ponderación. \\
\textbf{T} & Viernes 28 de agosto de 2026. Fase PPDIOO: \emph{Prepare}. \\
\bottomrule
\end{longtable}

\subsubsection{OE2 --- Especificar el protocolo de comunicación sensor--plataforma}

\begin{cajadefinicion}
\textbf{Especificar} un protocolo de comunicación entre sensor y plataforma basado en el estándar
Contact ID (SIA DC-05), \textbf{mediante} un análisis comparativo ponderado frente a alternativas
vigentes y la definición de un perfil de aplicación con su mapeo de códigos de evento, \textbf{para
obtener} un documento de especificación que incluya la gramática del mensaje, la tabla de equivalencias
y el contrato de la interfaz de ingesta, \textbf{al 25 de septiembre de 2026}.
\end{cajadefinicion}

\begin{longtable}{@{}L{1.5cm}L{12.2cm}@{}}
\filaencab \encab{} & \encab{Concreción} \\ \midrule \endhead
\textbf{S} & Especificación formal del protocolo y de la interfaz, no una descripción general. \\
\textbf{M} & $\geq 6$ alternativas evaluadas con matriz ponderada; $\geq 25$ códigos de evento mapeados; mensaje de evento $\leq 11$ bytes, verificado contra el payload útil mínimo de LoRaWAN; gramática formalizada. \\
\textbf{A} & El estándar SIA DC-05 es público y está documentado; el trabajo es de análisis y diseño. \\
\textbf{R} & Aborda directamente C1 y responde a la instrucción explícita del profesor guía de justificar la elección del protocolo frente a otros. \\
\textbf{T} & Viernes 25 de septiembre de 2026. Fase PPDIOO: \emph{Plan} y \emph{Design}. \\
\bottomrule
\end{longtable}

\subsubsection{OE3 --- Diseñar la arquitectura de la plataforma y su red de soporte}

\begin{cajadefinicion}
\textbf{Diseñar} la arquitectura de la plataforma y el diseño de red que la soporta, \textbf{mediante}
la aplicación de la fase \emph{Design} de PPDIOO y el uso de notaciones de modelado reconocidas,
\textbf{para obtener} un conjunto de diagramas, un modelo de datos y una especificación de seguridad y
segmentación que cubra la totalidad de los requerimientos levantados, \textbf{al 25 de septiembre de
2026}.
\end{cajadefinicion}

\begin{longtable}{@{}L{1.5cm}L{12.2cm}@{}}
\filaencab \encab{} & \encab{Concreción} \\ \midrule \endhead
\textbf{S} & Diagramas de arquitectura, topología de red, modelo de datos y diseño de seguridad. \\
\textbf{M} & $\geq 5$ diagramas en notación estándar; \textbf{100\%} de los requerimientos asignados a un componente; modelo de datos con $\geq 8$ entidades. \\
\textbf{A} & Trabajo documental de diseño, sin dependencia de adquisiciones. \\
\textbf{R} & Aborda C2 y C5; constituye el 55\% de la Etapa 2. \\
\textbf{T} & Viernes 25 de septiembre de 2026. Fase PPDIOO: \emph{Design}. \\
\bottomrule
\end{longtable}

\subsubsection{OE4 --- Determinar el modelo de despliegue y el costo de operación}

\begin{cajadefinicion}
\textbf{Determinar} la alternativa de alojamiento de la plataforma y su costo total de operación,
\textbf{mediante} la comparación de al menos tres opciones de despliegue bajo criterios de costo,
soberanía del dato y sostenibilidad, \textbf{para obtener} una recomendación justificada con el costo
mensual y anual estimado del producto, \textbf{al 25 de septiembre de 2026}.
\end{cajadefinicion}

\begin{longtable}{@{}L{1.5cm}L{12.2cm}@{}}
\filaencab \encab{} & \encab{Concreción} \\ \midrule \endhead
\textbf{S} & Recomendación única de despliegue con costo cuantificado. \\
\textbf{M} & $\geq 3$ alternativas comparadas; costo mensual y anual en pesos chilenos y dólares; matriz de criterios ponderados. \\
\textbf{A} & Información de precios de acceso público. \\
\textbf{R} & Responde a la instrucción del profesor guía de definir dónde se alojará la plataforma y cuánto costará el producto, aun siendo abierto. Aborda C3. \\
\textbf{T} & Viernes 25 de septiembre de 2026. Fase PPDIOO: \emph{Plan}. \\
\bottomrule
\end{longtable}

\subsubsection{OE5 --- Implementar la plataforma y sus adaptadores multiprotocolo}

\begin{cajadefinicion}
\textbf{Implementar} la interfaz de ingesta, el motor de normalización y la interfaz web de gestión,
junto con adaptadores para al menos dos transportes distintos, \textbf{mediante} desarrollo incremental
guiado por el backlog priorizado y versionado en repositorio git, \textbf{para obtener} un prototipo
funcional capaz de recibir, normalizar, almacenar y presentar datos de sensores de tecnologías
diferentes, \textbf{al 20 de noviembre de 2026}, con el 50\% de avance verificable al 23 de octubre.
\end{cajadefinicion}

\begin{longtable}{@{}L{1.5cm}L{12.2cm}@{}}
\filaencab \encab{} & \encab{Concreción} \\ \midrule \endhead
\textbf{S} & Prototipo funcional con ingesta, normalización, persistencia e interfaz de gestión. \\
\textbf{M} & \textbf{50\%} de las tareas de implementación al 23/10 y \textbf{100\%} al 20/11; $\geq 2$ transportes operativos; $\geq 1$ sensor físico reportando de extremo a extremo; mensaje descriptivo en el 100\% de los commits. \\
\textbf{A} & Condicionado a la disponibilidad de hardware; mitigado mediante emulación por software (riesgo R1). \\
\textbf{R} & Materializa la solución al problema central. \\
\textbf{T} & 23 de octubre (50\%) y 20 de noviembre de 2026 (100\%). Fase PPDIOO: \emph{Implement}. \\
\bottomrule
\end{longtable}

\subsubsection{OE6 --- Validar la plataforma en condiciones que replican el entorno operativo}

\begin{cajadefinicion}
\textbf{Validar} el funcionamiento de la plataforma en condiciones que replican el entorno operativo
real, \textbf{mediante} una batería de pruebas funcionales, de carga y de recuperación ante fallas, con
comparación cuantitativa antes y después de las optimizaciones aplicadas, \textbf{para obtener} un
informe de validación que acredite el cumplimiento de los requerimientos establecidos en OE1, \textbf{al
20 de noviembre de 2026}.
\end{cajadefinicion}

\begin{longtable}{@{}L{1.5cm}L{12.2cm}@{}}
\filaencab \encab{} & \encab{Concreción} \\ \midrule \endhead
\textbf{S} & Informe de validación con datos cuantitativos y trazabilidad a requerimientos. \\
\textbf{M} & $\geq 15$ casos de prueba documentados; $\geq$ 90\% aprobados; ingesta sostenida de $\geq 50$ mensajes por segundo sin pérdida; detección de sensor caído en $\leq 2$ intervalos de reporte; $\geq 3$ comparaciones cuantitativas antes y después de optimizar. \\
\textbf{A} & Las pruebas de carga se ejecutan con generación de tráfico sintético, sin requerir sensores físicos masivos. \\
\textbf{R} & Aborda C5 y E4; constituye el 30\% de la Etapa 3 y el 50\% de la Etapa 4. \\
\textbf{T} & Viernes 20 de noviembre de 2026. Fase PPDIOO: \emph{Operate} y \emph{Optimize}. \\
\bottomrule
\end{longtable}

\subsection{Matriz de trazabilidad}

\begin{longtable}{@{}L{1.3cm}L{3.9cm}L{2.1cm}L{1.6cm}L{2.4cm}C{1.1cm}@{}}
\filaencab \encab{Obj.} & \encab{Enunciado abreviado} & \encab{Causas} & \encab{Efectos} & \encab{Fase PPDIOO} & \encab{Etapa} \\ \midrule \endhead
\textbf{OE1} & Caracterizar y levantar requerimientos & C4 (diagnostica C1--C5) & E1, E3 & Prepare & 1 \\
\textbf{OE2} & Especificar el protocolo & \textbf{C1} & E2, E5 & Plan, Design & 2 \\
\textbf{OE3} & Diseñar arquitectura y red & C2, C5 & E2, E4 & Design & 2 \\
\textbf{OE4} & Despliegue y costo de operación & \textbf{C3} & E1, E5 & Plan & 2 \\
\textbf{OE5} & Implementar la plataforma & C1, C2, C4 & E1, E2, E3 & Implement & 3, 4 \\
\textbf{OE6} & Validar en entorno realista & C5 & E4 & Operate, Optimize & 4 \\
\bottomrule
\caption{Trazabilidad entre objetivos específicos, causas, efectos, fase metodológica y etapa.}
\label{tab:trazabilidad}
\end{longtable}

\subsubsection*{Verificación de cobertura}

\begin{longtable}{@{}L{5.6cm}L{2.5cm}L{4.4cm}L{1.6cm}@{}}
\filaencab \encab{Causa} & \encab{Objetivos} & \encab{Efecto} & \encab{Objetivos} \\ \midrule \endhead
C1 Fragmentación de protocolos & OE1, \textbf{OE2}, OE5 & E1 Inversión inmovilizada & OE1, OE4, OE5 \\
C2 Acoplamiento hardware-plataforma & \textbf{OE3}, OE5 & E2 Datos perdidos & OE2, OE3, OE5 \\
C3 Costo y dependencia comercial & \textbf{OE4} & E3 Decisiones sin dato & OE1, OE5 \\
C4 Falta de capacidad técnica & \textbf{OE1}, OE5 & E4 Fallas detectadas tarde & OE3, \textbf{OE6} \\
C5 Sin visibilidad de conectividad & OE3, \textbf{OE6} & E5 Dependencia de proveedor & OE2, OE4 \\
\bottomrule
\caption{Cobertura de causas y efectos por los objetivos específicos.}
\label{tab:cobertura}
\end{longtable}

Las cinco causas y los cinco efectos quedan cubiertos por al menos un objetivo específico. Ninguna causa
carece de objetivo asociado y ningún objetivo carece de vínculo con el problema.
